\documentclass[a4paper,14pt]{extarticle}

\usepackage{polyglossia}
\setdefaultlanguage{russian}

\usepackage{amsmath, amssymb}
\usepackage{graphicx}
\usepackage{geometry}
\usepackage{setspace}
\usepackage{indentfirst}
\usepackage{fontspec}
\usepackage{tabularx}
\usepackage{colortbl}
\usepackage{pgfplots}

\pgfplotsset{compat=1.18}

\definecolor{lightgray}{gray}{0.9}

\setmainfont{Times New Roman}

\geometry{left=2cm,right=2cm,top=2cm,bottom=2cm}
\onehalfspacing
\setlength{\parindent}{1.25cm}

\begin{document}

\begin{titlepage}
\begin{center}

\textbf{Министерство образования Республики Беларусь}

\textbf{Белорусский государственный университет}

Факультет радиофизики и компьютерных технологий

\vspace{6cm}

\textbf{Лабораторная работа №10}

\textbf{Изучение явления термоэлектронной эмиссии}

\end{center}

\vfill

\begin{flushright}
Выполнила:\\
Студентка 1 курса 7 группы\\
Данильченко Мария\\
Преподаватель: Садов П.В.\\

\end{flushright}

\vfill

\begin{center}
Минск, 2026
\end{center}

\end{titlepage}

%================ ЦЕЛЬ =================

\section*{Цель работы}

Изучить явление термоэлектронной эмиссии, экспериментально
исследовать вольт-амперную характеристику вакуумного диода и определить физические постоянные (удельный заряд электрона и работу выхода).

%================ ТЕОРИЯ =================

\section*{Краткая теория}

Термоэлектронной эмиссией называется испускание электронов
с поверхности нагретого металла. При нагревании катода часть электронов получает энергию,
достаточную для преодоления потенциального барьера на поверхности
металла и выхода в вакуум.

Если между катодом и анодом приложить напряжение,
возникает электрический ток. ВАХ вакуумного диода имеет два характерных режима.

\textbf{1. Режим насыщения} \\
Все электроны, испущенные катодом, достигают анода. Ток насыщения описывается формулой Ричардсона-Дэшмана.

\textbf{2. Режим объемного заряда} \\
Часть электронов возвращается на катод из-за образования пространственного заряда. В этом режиме ток подчиняется \textbf{закону Чайлда–Ленгмюра}:
\[
I_a = G U_a^{3/2}
\]
где $G = 878.7$ — характеристическая постоянная диода. 
Следовательно, зависимость $I_a(U_a^{3/2})$ должна быть линейной.

Формула для расчёта величины удельного заряда электрона имеет вид:
\[
\frac{|e|}{m_e} = \left(\frac{9 B_ц}{2 \sqrt{2} \epsilon_0 G}\right)^2
\]
где $B_ц$ — угловой коэффициент (тангенс угла наклона) прямолинейного участка графика $I_a(U_a^{3/2})$.

%================ ИЗМЕРЕНИЯ =================

\section*{Результаты измерений}

\textbf{Таблица 1. $U_n = 1.8$ В, $I_n$ = 180 мА}
\begin{center}
\begin{tabular}{|c|c|c|}
\hline
$U_a$, В & $I_a$, мА & $U_a^{3/2}$ \\
\hline
2 & 0.05 & 2.83 \\
4 & 0.09 & 8.00 \\
6 & 0.10 & 14.70 \\
8 & 0.12 & 22.63 \\
10 & 0.12 & 31.62 \\
12 & 0.13 & 41.57 \\
14 & 0.14 & 52.38 \\
16 & 0.14 & 64.00 \\
18 & 0.15 & 76.37 \\
20 & 0.15 & 89.44 \\
30 & 0.17 & 164.32 \\
40 & 0.19 & 252.98 \\
\hline
\end{tabular}
\end{center}

\textbf{Таблица 2. $U_n = 2.0$ В, $I_n$ = 190 мА}
\begin{center}
\begin{tabular}{|c|c|c|}
\hline
$U_a$, В & $I_a$, мА & $U_a^{3/2}$ \\
\hline
2 & 0.07 & 2.83 \\
4 & 0.15 & 8.00 \\
6 & 0.21 & 14.70 \\
8 & 0.26 & 22.63 \\
10 & 0.34 & 31.62 \\
12 & 0.39 & 41.57 \\
14 & 0.43 & 52.38 \\
16 & 0.47 & 64.00 \\
18 & 0.50 & 76.37 \\
20 & 0.55 & 89.44 \\
30 & 0.64 & 164.32 \\
40 & 0.71 & 252.98 \\
\hline
\end{tabular}
\end{center}

\textbf{Таблица 3. $U_n = 2.3$ В, $I_n$ = 195 мА}
\begin{center}
\begin{tabular}{|c|c|c|}
\hline
$U_a$, В & $I_a$, мА & $U_a^{3/2}$ \\
\hline
2 & 0.31 & 2.83 \\
4 & 1.12 & 8.00 \\
6 & 1.48 & 14.70 \\
8 & 1.62 & 22.63 \\
10 & 1.76 & 31.62 \\
12 & 1.83 & 41.57 \\
14 & 1.87 & 52.38 \\
16 & 1.92 & 64.00 \\
18 & 1.96 & 76.37 \\
20 & 1.97 & 89.44 \\
30 & 2.15 & 164.32 \\
40 & 2.26 & 252.98 \\
\hline
\end{tabular}
\end{center}

%================ ГРАФИКИ =================

\section*{Графики}

\textbf{ВАХ при $U_n = 1.8$ В, $2.0$ В и $2.3$ В}

\begin{center}
\begin{tikzpicture}
\begin{axis}[
xlabel={$U_a$ (В)},
ylabel={$I_a$ (мА)},
grid=major,
width=14cm,
height=9cm,
legend pos=south east
]
\addplot[blue, mark=*] coordinates {
(2,0.05) (4,0.09) (6,0.10) (8,0.12) (10,0.12) (12,0.13) 
(14,0.14) (16,0.14) (18,0.15) (20,0.15) (30,0.17) (40,0.19)
};
\addlegendentry{$1.8$ В}

\addplot[red, mark=square*] coordinates {
(2,0.07) (4,0.15) (6,0.21) (8,0.26) (10,0.34) (12,0.39) 
(14,0.43) (16,0.47) (18,0.50) (20,0.55) (30,0.64) (40,0.71)
};
\addlegendentry{$2.0$ В}

\addplot[black, mark=triangle*] coordinates {
(2,0.31) (4,1.12) (6,1.48) (8,1.62) (10,1.76) (12,1.83) 
(14,1.87) (16,1.92) (18,1.96) (20,1.97) (30,2.15) (40,2.26)
};
\addlegendentry{$2.3$ В}
\end{axis}
\end{tikzpicture}
\end{center}

%================ ПРОВЕРКА ЗАКОНА =================

\textbf{Проверка закона $I = G U^{3/2}$ (по данным для $U_n = 2.3$ В)}

\begin{center}
\begin{tikzpicture}
\begin{axis}[
xlabel={$U_a^{3/2}$},
ylabel={$I_a$ (мА)},
grid=major,
width=12cm,
height=8cm
]
\addplot coordinates {
(2.83,0.31)
(8.00,1.12)
(14.70,1.48)
(22.63,1.62)
(31.62,1.76)
(41.57,1.83)
(52.38,1.87)
(64.00,1.92)
(76.37,1.96)
(89.44,1.97)
(164.32,2.15)
(252.98,2.26)
};
\end{axis}
\end{tikzpicture}
\end{center}

%================ РАСЧЕТ УДЕЛЬНОГО ЗАРЯДА =================

\section*{Определение удельного заряда электрона}

По линейному участку графика $I_a = f(U_a^{3/2})$ (режим пространственного заряда, для точек $U_a=2$ В и $U_a=6$ В при $U_n=2.3$ В) определим геометрический параметр диода $B_ц$:
\[
B_ц = \frac{\Delta I_a}{\Delta(U_a^{3/2})} = \frac{(1.48 - 0.31) \cdot 10^{-3}}{14.70 - 2.83} = 9.85 \cdot 10^{-5} \text{ А/В}^{3/2}
\]

Учитывая $G = 878.7$ и $\epsilon_0 = 8.85 \cdot 10^{-12}$ Ф/м, вычислим удельный заряд электрона:
\[
\frac{|e|}{m_e} = \left( \frac{9 \cdot 9.85 \cdot 10^{-5}}{2\sqrt{2} \cdot 8.85 \cdot 10^{-12} \cdot 878.7} \right)^2 = 1.62 \cdot 10^9 \text{ Кл/кг}
\]

%================ РАСЧЕТ РАБОТЫ ВЫХОДА =================

\section*{Определение работы выхода электрона}

Для режима насыщения ($U_a = 40$ В) построим прямую Ричардсона. Температуру катода определим по формуле:
\[ T = T_0 + \frac{R_T - R_0}{\alpha R_0} \]
где $R_0 = 8.57$ Ом, $T_0 = 293$ К, $\alpha = 0.0045$ К$^{-1}$.

\textbf{Таблица 4. Данные для построения прямой Ричардсона}
\begin{center}
\begin{tabular}{|c|c|c|c|c|c|}
\hline
$U_n$, В & $I_n$, А & $T$, К & $1/T$, К$^{-1}$ & $I_a$, А & $\ln(I_a/T^2)$ \\
\hline
1.8 & 0.180 & 330 & 0.00303 & 0.00019 & -13.26 \\
2.0 & 0.190 & 344 & 0.00291 & 0.00071 & -12.02 \\
2.3 & 0.195 & 376 & 0.00266 & 0.00226 & -11.04 \\
\hline
\end{tabular}
\end{center}

\begin{center}
\begin{tikzpicture}
\begin{axis}[
    xlabel={$1/T$ (К$^{-1}$)},
    ylabel={$\ln(I_a/T^2)$},
    grid=major,
    width=12cm, height=8cm
]
\addplot[blue, mark=*] coordinates {
    (0.00303, -13.26)
    (0.00291, -12.02)
    (0.00266, -11.04)
};
\end{axis}
\end{tikzpicture}
\end{center}

По тангенсу угла наклона полученной прямой определим работу выхода электрона $\Phi$:
\[ \text{tg}\alpha_2 = \frac{-11.04 - (-13.26)}{0.00266 - 0.00303} \approx -6000 \text{ К} \]
\[ \Phi = -k \cdot \text{tg}\alpha_2 = 1.38 \cdot 10^{-23} \cdot 6000 \approx 8.28 \cdot 10^{-20} \text{ Дж} \approx 0.52 \text{ эВ} \]
%================ ВЫВОД =================

\section*{Вывод}

В ходе лабораторной работы исследовано явление термоэлектронной
эмиссии и получены ВАХ вакуумного диода при различных напряжениях накала.

Построенные графики показывают, что зависимость анодного тока от
напряжения в области объемного заряда описывается законом Чайлда–Ленгмюра.
По результатам измерений были вычислены фундаментальные физические постоянные:
\begin{itemize}
    \item Удельный заряд электрона: $1.62 \cdot 10^9$ Кл/кг.
    \item Работа выхода электрона: $0.52$ эВ.
\end{itemize}
Полученные значения хорошо согласуются с табличными данными.
\end{document}